\subsection{トピックと参考文献の対応}

この表は、SWEBOKが示す「トピック対参考資料の行列」を、学習用に簡略化したものである。詳しく学ぶときは、各トピックに対応する文献を読む。

\par\smallskip\noindent\refstepcounter{table}\label{tab:ch03-08}表 \thetable: トピックと参考文献の対応の整理\par\smallskip
表 \thetable は、トピックと参考文献の対応の整理を比較・確認しやすい形で整理したものである。\par\smallskip
\begin{longtable}{L{66mm}L{98mm}}
\toprule
トピック & 主な参考資料 \\
\midrule
\endhead
1. ソフトウェア設計の基礎 & Brooks, The Design of Design; Budgen, Software Design \\
1.1 設計思考 & Brooks; Budgen; Ross, Goodenough and Irvine \\
1.2 ソフトウェア設計の文脈 & Budgen; Sommerville \\
1.3 ソフトウェア設計の主要論点 & Brooks; Budgen; Sommerville; Bosch \\
1.4 ソフトウェア設計原則 & Budgen; Dijkstra; Parnas; ISO/IEC/IEEE 24765; IEEE Std 7000 \\
2. ソフトウェア設計プロセス & Budgen; Sommerville; ISO/IEC/IEEE 12207 \\
2.1 高水準設計 & Brooks; Budgen \\
2.2 詳細設計 & Budgen; Sommerville \\
3. ソフトウェア設計品質 & Budgen; Sommerville; Galster et al. \\
3.1 並行性 & Sommerville \\
3.2 制御とイベント処理 & Sommerville; Luckham; Mühl, Fiege and Pietzuch \\
3.3 データ永続化 & Sommerville \\
3.4 構成要素の分散 & Sommerville \\
3.5 誤り・例外処理と耐故障性 & Sommerville \\
3.6 統合と相互運用性 & Budgen \\
3.7 保証・セキュリティ・安全性 & Sommerville \\
3.8 変動性 & Galster et al.; Bosch \\
4. ソフトウェア設計の記録 & Booch, Rumbaugh and Jacobson; Budgen; Sommerville \\
4.1 モデルに基づく設計 & Budgen; Sommerville \\
4.2 構造設計記述 & Booch et al.; Budgen; Gamma et al.; Sommerville \\
4.3 振る舞い設計記述 & Booch et al.; Budgen; Gamma et al.; Sommerville \\
4.4 設計パターンとスタイル & Brooks; Budgen; Gamma et al.; Sommerville \\
4.5 専用言語と領域固有言語 & Sommerville; Kosar, Bohra and Mernik \\
4.6 設計根拠 & Brooks; Budgen; Sommerville; Parnas and Clements \\
5. ソフトウェア設計戦略と方法 & Sommerville; Budgen; Nielsen; Kiczales et al. \\
5.1 一般戦略 & Budgen \\
5.2 機能指向・構造化設計 & Budgen \\
5.3 データ中心設計 & Budgen; Sommerville \\
5.4 オブジェクト指向設計 & Budgen; Booch et al. \\
5.5 利用者中心設計 & Brooks; Nielsen \\
5.6 構成要素ベース設計 & Booch et al.; Budgen; Sommerville \\
5.7 イベント駆動設計 & Sommerville; Luckham; Mühl et al. \\
5.8 アスペクト指向設計 & Booch et al.; Sommerville; Kiczales et al. \\
5.9 制約に基づく設計 & Brooks \\
5.10 ドメイン駆動設計 & Budgen \\
5.11 その他の方法 & Sommerville \\
6. 設計品質分析と評価 & Budgen; Sommerville; Parnas and Weiss \\
6.1 設計レビューと監査 & Budgen; Parnas and Weiss \\
6.2 品質属性 & Sommerville \\
6.3 品質分析・評価技法 & Sommerville; Software Quality KA \\
6.4 測度とメトリクス & Budgen; Sommerville \\
6.5 検証・妥当性確認・認証 & Sommerville; Software Quality KA \\
\bottomrule
\end{longtable}
